\documentclass[../libro]{subfiles}

\begin{document}

\ifSubfilesClassLoaded{\appendix\chapter{后日谈}\clearpage}{}

\section{Until the re-seeing}

\begingroup
\setlength{\parindent}{1.5em}%for 12pt

The title of this section is a literal translation of
``Ĝis la revido'' in Esperanto,
which idiomatically means ``See you'' in English.
Yes, I am bidding farewell---it is not the case that
\textit{everything} that has a beginning has an end,
but it \textit{is} the case that
this book,
\textit{An introduction to determinants},
which has a beginning, has an end.

I do not really have anything interesting to say,
but I still choose to say something here.

This book is mainly about determinants,
and it is over four hundred pages long.
I know that I have written too much in this book,
but I really wanted to share my knowledge,
hoping that it might be useful and helpful.

This book is available under
the \href{https://opensource.org/license/0bsd}{Zero-Clause BSD License},
which means that
any person or organisation is allowed to use this book
in any way, for any purpose, with or without charge.
I have not yet earned a single cent from this book,
but I did not, do not, and will not
prevent any person or organisation
from making a profit from this book.
I want to propagate the values of ``open knowledge'',
and I want to prove its profitability,
if I can ever earn money from this book.

I thank mathematicians and their work;
I thank artists who made artworks for this book;
I thank my friends, who were good brokers;
I thank everyone who read or reads this book.

Parenthetically, the day on which
I wrote this closing remark is 28 June 2025.
Both six and twenty-eight are perfect numbers,
in the sense that
each is equal to the sum of all its positive divisors,
excluding itself.
It is also worth noting that
\(6.28\) is a decimal approximation of \(2\pi\),
the ratio of the circumference of a circle to its radius.
Therefore I think of 28 June as
a ``mirinda'' (meaning ``worth marvelling at'') day.

Fartu bone. Kredu, ke ni renkontiĝos iam, ie, iel, ial.
(Do well. Believe that we will meet
sometime, somewhere, somehow, for some reason.)

\endgroup

\end{document}

This is just a closing remark.
